%% =========================================================================
%%  mvher: Missing-Value Handling for Entity Resolution in R
%%  Submitted to The R Journal
%%  Author: Xinmin Chu
%% =========================================================================
\documentclass[a4paper]{article}
\usepackage[utf8]{inputenc}
\usepackage[T1]{fontenc}
\usepackage{RJournal}
\usepackage{amsmath,amssymb,array}
\usepackage{booktabs}
\usepackage{microtype}
\usepackage{graphicx}
\usepackage{natbib}

\begin{document}

\sectionhead{Contributed research article}
\RJvolume{XX}
\RJvolnumber{X}
\RJmonth{XXXXX}

\begin{article}

%%  ── Title and author ───────────────────────────────────────────────────
\title{\pkg{mvher}: Missing-Value Handling for Entity Resolution in R}

\author{Xinmin Chu \and David Degras-Valabregue}
\date{}

\maketitle

%%  ── Abstract ───────────────────────────────────────────────────────────
\RJabstract{%
Missing values are common in entity resolution (ER) and record
linkage applications, yet most pipelines handle them inconsistently---
either ignoring them or applying generic imputation rules designed for
non-relational settings. A particularly subtle artifact arises from
standard string-similarity metrics such as Jaro-Winkler and Levenshtein:
two empty strings are scored as \emph{perfectly} similar (similarity~$=
1$), sending a spurious ``definite match'' signal to any classifier that
consumes pairwise features. Under entity-disjoint train/test
evaluation---the methodologically correct protocol for ER---this
artifact causes systematic covariate shift when missingness is
concentrated within certain entities. We introduce \pkg{mvher}, an R
package that addresses these problems through a principled, ER-specific
framework. The package provides (i)~four diagnostic functions for
measuring missingness rates, detecting the underlying mechanism (MCAR,
MAR, or MNAR), and quantifying entity-level intra-class correlation of
missingness; (ii)~eight field-level comparison strategies that handle
empty and \texttt{NA} values explicitly during pairwise similarity
computation; (iii)~seven feature-matrix imputation strategies with a
fit-on-training/apply-to-test API that prevents data leakage; and
(iv)~an automated recommendation engine that selects the best strategy
combination from data characteristics. We illustrate the package on
synthetic and real-world benchmarks and show that the \texttt{"full"}
comparison strategy---neutral similarity plus binary missingness
indicators---substantially reduces bias in link probability estimates
compared to the naive both-null-equals-one default.
\pkg{mvher} is available at \url{https://github.com/xinminchu/mvher}.
}

%%  ── 1. Introduction ────────────────────────────────────────────────────
\section{Introduction}

Entity resolution (ER)---also known as \emph{record linkage}
\citep{Newcombe1959}, \emph{deduplication} \citep{Christen2012}, or
\emph{entity matching} \citep{Mudgal2018}---is the task of identifying
records across one or more databases that refer to the same real-world
entity. Applications span patient record integration in clinical
systems, bibliographic deduplication, business register linkage, and
product catalog consolidation \citep{Christophides2020}.

A standard supervised ER pipeline has three stages.  First,
\emph{blocking} restricts pairwise comparisons to candidate pairs within
homogeneous ``blocks,'' reducing the quadratic cost to a manageable
number of comparisons \citep{Christen2012}.  Second, for each candidate
pair, \emph{field comparison} computes pairwise string similarities
across shared fields (name, address, date of birth, etc.) to produce a
feature vector.  Third, a \emph{classifier} scores each candidate pair
as a match or non-match; the resulting decisions are grouped into
entity clusters.

Missing values enter this pipeline at the field-comparison stage.
Administrative records, scraped web data, and integrated
multi-source tables routinely contain fields that are empty or
\texttt{NA} for a significant fraction of records.  Despite this, the
treatment of missing values is rarely discussed in ER methodological
papers, and the default behavior of popular R packages is to score two
empty strings as perfectly similar (Jaro-Winkler or Levenshtein
similarity~$= 1$).

This default produces two distinct failure modes.

\begin{enumerate}
  \item \textbf{Signal corruption.} A pair of records that are both
    missing a field receives a similarity of~$1$ for that field---the
    same score as two identical non-empty strings. Any classifier that
    uses this feature will learn a spuriously large positive weight for
    the field, inflating estimated match probabilities for pairs that
    share a missing value.

  \item \textbf{Covariate shift under entity-disjoint splits.}
    Entity-disjoint evaluation---where training and test entities are
    disjoint---is the correct protocol for ER, because it mirrors
    deployment conditions in which the classifier encounters entirely
    new entities \citep{Murray2016, Christen2012}. When missingness is
    entity-correlated (e.g., one data source systematically omits a
    field for certain entities), entities that are missing in the test
    set were not seen during training. The spurious both-null-$= 1$
    signal then has a different marginal distribution at test time than
    at train time, creating covariate shift that degrades classifier
    performance in a way that is invisible under record-disjoint
    evaluation.
\end{enumerate}

There is a rich literature on handling missing data in general
statistical settings \citep{Rubin1976, LittleRubin2002, vanBuuren2018},
but very little guidance specific to the pairwise-comparison structure
of ER pipelines. The field-level strategies developed for the
Fellegi-Sunter model \citep{FellegiSunter1969} encode a three-level
comparison ($\gamma = 0$: no agreement, $\gamma = 1$: partial agreement,
$\gamma = 2$: full agreement) that naturally accommodates missing fields,
but this encoding is rarely implemented correctly in practice and does
not integrate cleanly with modern machine-learning classifiers that
expect continuous features.

We present \pkg{mvher}, an R package that fills this gap. The package is
designed around four principles:

\begin{enumerate}
  \item \textbf{Transparency.} Every strategic choice made about missing
    values is made explicit and documented, not buried inside a helper
    function.
  \item \textbf{Leakage prevention.} Imputation parameters are always
    estimated on training data and applied without re-estimation to test
    data, preventing the subtle leakage that arises when test-set
    statistics inform imputed training values.
  \item \textbf{ER-native strategies.} The package includes comparison
    strategies---neutral similarity, asymmetric handling, Fellegi-Sunter
    encoding, and binary missingness indicators---that are specifically
    motivated by the pairwise structure of ER feature extraction, not
    by general tabular imputation.
  \item \textbf{Guided selection.} An automated recommendation engine
    maps data characteristics (missingness rate, entity-level
    correlation, split type) to the appropriate strategy combination,
    reducing the burden on practitioners who may be unfamiliar with the
    theoretical foundations.
\end{enumerate}

The rest of this paper is organized as follows. Section~\ref{sec:bg}
reviews ER pipelines, missingness mechanisms, and the
both-null artifact that motivates the package. Section~\ref{sec:pkg}
provides an overview of the \pkg{mvher} package design and API.
Sections~\ref{sec:diag}--\ref{sec:recommend} walk through each module
in detail with reproducible R code.  Section~\ref{sec:casestudy}
presents an end-to-end case study illustrating the full workflow.
Section~\ref{sec:compare} positions \pkg{mvher} relative to existing
packages.  Section~\ref{sec:summary} summarizes contributions and
outlines directions for future work.

%%  ── 2. Background ──────────────────────────────────────────────────────
\section{Background}
\label{sec:bg}

\subsection{Entity resolution pipelines}

The supervised ER pipeline produces a feature matrix $\mathbf{X}$ whose
row $k$ corresponds to candidate pair $(i, j)$, and whose columns
are pairwise similarity scores $s_f(r_i, r_j)$ for each field $f$.
Under the Fellegi-Sunter model \citep{FellegiSunter1969}, records are
classified by comparing the likelihood ratio of the comparison vector
$\boldsymbol{\gamma}$ under the match and non-match populations.
Modern supervised alternatives train logistic regression, random
forests, gradient boosting, or neural network classifiers on labeled
pairs \citep{Christophides2020, Mudgal2018}.

Evaluation requires held-out data.  We focus on
\emph{entity-disjoint} evaluation, where all records belonging to a
random subset of entities are reserved as the test set
\citep{Murray2016}.  Under this protocol, when missingness is
concentrated within certain entities, the training and test feature
distributions can differ systematically---a problem we return to in
Section~\ref{sec:bg:artifact}.

\subsection{Missingness mechanisms}

\citet{Rubin1976} classified missing-data mechanisms into three types,
a taxonomy that directly informs the choice of handling strategy.

\begin{description}
  \item[MCAR (Missing Completely At Random).] The probability that a
    value is missing is independent of any data values, observed or
    unobserved.  Listwise deletion and mean imputation are consistent
    under MCAR.

  \item[MAR (Missing At Random).] The probability of missingness can
    depend on observed values but not on the missing value itself.
    Most multiple imputation methods are valid under MAR.

  \item[MNAR (Missing Not At Random).] The probability of missingness
    depends on the unobserved value.  No imputation method can
    fully recover from MNAR without additional assumptions.
\end{description}

In ER settings, MNAR is common: a field is often absent because the
underlying value is unknown, not because it was randomly omitted.
More critically, missingness is frequently \emph{entity-correlated}:
an entire record from a particular data source may lack a field because
that source does not collect it.  This entity-level clustering of
missingness can be quantified by the intra-class correlation coefficient
(ICC) of the missingness indicator with respect to entity membership
\citep{ShroutFleiss1979}.  When the ICC is high and the evaluation uses
entity-disjoint splits, the field-comparison features in training and
test sets have systematically different marginal distributions, creating
covariate shift.

\subsection{The both-null artifact}
\label{sec:bg:artifact}

Let $d(s, t)$ denote a string edit distance (e.g., normalized
Levenshtein), and let $\mathrm{sim}(s, t) = 1 - d(s, t)$ be the
corresponding similarity.  For any standard metric, $d(\epsilon,
\epsilon) = 0$, so $\mathrm{sim}(\epsilon, \epsilon) = 1$ for the empty
string~$\epsilon$.  Similarly, $\mathrm{JW}(\epsilon, \epsilon) = 1$
under the Jaro-Winkler definition \citep{Jaro1989, Winkler1990}.

When both records in a pair are missing a field, this convention
assigns a similarity of~1---identical to the score two records would
receive if they shared the same non-empty value for that field. A logistic
regression or other linear classifier that receives this signal will
learn a positive weight for the field, associating high similarity with
match status. Pairs that are both missing the field will therefore
receive an inflated predicted match probability, even if they are not
the same entity.

This problem is made worse when, under entity-disjoint splits, the
both-missing rate is substantially higher at test time than at training
time (because the test entities have a higher base rate of missingness
for that field). The classifier's decision boundary was calibrated on
training data where the spurious signal appeared less often; it will be
systematically poorly calibrated on test data.

%%  ── 3. Package overview ────────────────────────────────────────────────
\section{Package overview}
\label{sec:pkg}

\pkg{mvher} is available on GitHub at
\url{https://github.com/xinminchu/mvher} and can be installed with:

\begin{example}
# install.packages("remotes")
remotes::install_github("xinminchu/mvher")
\end{example}

The package requires R~$\geq 4.0.0$ and depends on
\CRANpkg{data.table} \citep{datatable} for memory-efficient data
manipulation and \CRANpkg{stringdist} \citep{stringdist} for string
similarity computation.  It exports 13~functions organized into four
modules:

\begin{description}
  \item[\textbf{Diagnosis}] (\texttt{mv\_field\_summary},
    \texttt{mv\_pair\_summary}, \texttt{mv\_detect\_mechanism},
    \texttt{mv\_entity\_correlation}): characterize the extent, pattern,
    and mechanism of missingness in a record set.

  \item[\textbf{Field comparison}] (\texttt{mv\_compare\_fields},
    \texttt{mv\_pair\_features}): compute pairwise string similarity
    features with explicit control over missing-value behavior.

  \item[\textbf{Imputation}] (\texttt{mv\_fit\_imputer},
    \texttt{mv\_apply\_imputer}, \texttt{mv\_impute}): fill residual
    \texttt{NA}s in feature matrices with a fit-once/apply-many API.

  \item[\textbf{Strategy guidance}] (\texttt{mv\_strategy\_table},
    \texttt{mv\_recommend}): enumerate available strategies and
    automatically select the best combination.
\end{description}

A utility function \texttt{norm\_str()} performs standard text
normalization (lowercasing, trimming, whitespace collapsing,
\texttt{NA} coercion to empty string) and is exported for use in user
preprocessing pipelines.

The four modules are composed sequentially in a typical pipeline:
diagnose $\to$ compare fields $\to$ impute residual NAs $\to$ classify.

%%  ── 4. Diagnosing missingness ──────────────────────────────────────────
\section{Diagnosing missingness}
\label{sec:diag}

The first step in any principled missing-value strategy is to understand
the extent and structure of missingness in the data. \pkg{mvher}
provides four diagnostic functions for this purpose.

\subsection{Field-level and pair-level summaries}

\texttt{mv\_field\_summary()} returns per-column missingness rates,
counting both \texttt{NA} and empty strings (\texttt{""}) as missing:

\begin{example}
library(mvher)

records <- data.frame(
  id        = 1:6,
  entity_id = c(1L, 1L, 2L, 2L, 3L, 3L),
  name      = c("Alice Smith", "Alice Smyth", NA,
                "Bob Jones", "", "Carol White"),
  city      = c("London", "Londen", "Paris", NA, "Rome", "Roma"),
  stringsAsFactors = FALSE
)

mv_field_summary(records, c("name", "city"))
#>    column n_total n_empty  pct_miss
#> 1:   name       6       2 0.3333333
#> 2:   city       6       1 0.1666667
\end{example}

At the pair level, \texttt{mv\_pair\_summary()} reports how many
candidate pairs are affected by missingness in each field, broken down
by whether one or both records are missing:

\begin{example}
pairs <- data.frame(id1 = c(1, 2, 3, 5), id2 = c(2, 4, 4, 6))
mv_pair_summary(records, pairs, c("name", "city"))
#>    column n_pairs any_missing both_missing one_missing
#> 1:   name       4           3            1           2
#> 2:   city       4           1            0           1
\end{example}

This pair-level view is more directly informative than the record-level
summary: a high \texttt{both\_missing} count for a field indicates that
the both-null artifact will affect many candidate pairs.

\subsection{Mechanism detection}

\texttt{mv\_detect\_mechanism()} applies a chi-squared test of
independence between the missingness indicator for each field and entity
membership. The test is approximate---it treats the entity column as a
categorical factor and uses the Pearson $\chi^2$ statistic---but it
provides a useful first-pass signal:

\begin{example}
mv_detect_mechanism(records, c("name", "city"))
#>    column overall_miss_rate entity_miss_variance chi_sq_p mechanism
#> 1:   name         0.3333333            0.1666667 0.135335   unknown
#> 2:   city         0.1666667            0.1666667 0.135335   unknown
\end{example}

The \texttt{mechanism} column is classified as \texttt{"MCAR"} when the
chi-squared p-value exceeds the significance threshold $\alpha$ (default
$0.05$); \texttt{"MAR/MNAR"} when $p < \alpha$; and \texttt{"unknown"}
when the test cannot be performed (e.g., too few records per entity).

\subsection{Entity-level intra-class correlation}

The most diagnostic measure for the entity-disjoint split problem is the
intra-class correlation coefficient (ICC) of the missingness indicator
with respect to entity membership. \pkg{mvher} computes a one-way ICC
following \citet{ShroutFleiss1979}:

\begin{equation}
\mathrm{ICC} = \frac{\hat{\sigma}^2_{\text{between}}}
               {\hat{\sigma}^2_{\text{between}} + \hat{\sigma}^2_{\text{within}}},
\label{eq:icc}
\end{equation}

where $\hat{\sigma}^2_{\text{between}}$ and
$\hat{\sigma}^2_{\text{within}}$ are the between-entity and
within-entity variance components of a one-way ANOVA on the binary
missingness indicator, estimated by the method of moments.

\texttt{mv\_entity\_correlation()} returns the ICC and an interpretation
label:

\begin{example}
mv_entity_correlation(records, c("name", "city"))
#>    column  icc interpretation
#> 1:   name 0.50 high (MNAR-like)
#> 2:   city 0.25 moderate (MAR-like)
\end{example}

An ICC $> 0.3$ indicates that missingness is substantially clustered
within entities. When combined with entity-disjoint splits, this signals
high risk of covariate shift and recommends the use of the
\texttt{"full"} comparison strategy (Section~\ref{sec:field}).

%%  ── 5. Field-comparison strategies ────────────────────────────────────
\section{Field-comparison strategies}
\label{sec:field}

\texttt{mv\_compare\_fields(xa, xb, method, ...)} computes pairwise
similarities between two character vectors \texttt{xa} and \texttt{xb}
under one of eight strategies, specified by the \texttt{method}
argument. The function always returns Jaro-Winkler (\texttt{jw}) and
normalized Levenshtein (\texttt{lv}) similarities along with binary
indicators \texttt{any\_miss} and \texttt{both\_miss}, with additional
outputs depending on the strategy.

Table~\ref{tab:strategies} summarizes the eight strategies and their
key properties.

\begin{table}[ht]
\centering
\caption{Field-comparison strategies in \pkg{mvher}. ``MNAR-safe''
  means the strategy does not assign a spuriously high similarity to
  both-missing pairs. ``Adds features'' means extra indicator or
  encoding columns are appended to the feature vector.}
\label{tab:strategies}
\smallskip
\begin{tabular}{lp{4.8cm}cc}
\toprule
\textbf{Method} & \textbf{Behavior for missing fields}
  & \textbf{MNAR-safe} & \textbf{Adds features} \\
\midrule
\texttt{"current"}       & both-empty $\to$ sim$=1$; else raw sim  & No  & No \\
\texttt{"neutral"}       & any-empty $\to$ sim$=0.5$               & Yes & No \\
\texttt{"zero"}          & any-empty $\to$ sim$=0$                 & Yes & No \\
\texttt{"asymmetric"}    & both-empty $\to 0$; one-empty $\to 0.5$ & Yes & No \\
\texttt{"indicator"}     & raw sim $+$ \texttt{any\_miss}/\texttt{both\_miss} & Partial & Yes \\
\texttt{"full"}          & neutral sim $+$ \texttt{any\_miss}/\texttt{both\_miss} & Yes & Yes \\
\texttt{"fellegi\_sunter"} & integer $\gamma$-level encoding 0--4  & Yes & Yes \\
\texttt{"complete\_cases"} & any-empty $\to$ \texttt{NA}           & Yes & No \\
\bottomrule
\end{tabular}
\end{table}

The following example shows the behavior of five methods on a short
vector with one non-missing, one partially-missing, and one
both-missing pair:

\begin{example}
xa <- c("Alice Smith", "Bob Jones", "",     "")
xb <- c("Alice Smyth", "Robert J.", "Carol", "")

# Default: both-empty scores 1 (the both-null artifact)
mv_compare_fields(xa, xb, method = "current")$jw
#> [1] 0.9666667 0.8162963 0.0000000 1.0000000

# Neutral: both-empty -> 0.5, one-empty -> 0.5
mv_compare_fields(xa, xb, method = "neutral")$jw
#> [1] 0.9666667 0.8162963 0.5000000 0.5000000

# Zero: any missing -> 0
mv_compare_fields(xa, xb, method = "zero")$jw
#> [1] 0.9666667 0.8162963 0.0000000 0.0000000

# Asymmetric: both-empty -> 0, one-empty -> 0.5
mv_compare_fields(xa, xb, method = "asymmetric")$jw
#> [1] 0.9666667 0.8162963 0.5000000 0.0000000

# Fellegi-Sunter integer gamma levels (0 = absent, 4 = full agreement)
mv_compare_fields(xa, xb, method = "fellegi_sunter")$gamma
#> [1] 4 3 0 0
\end{example}

\subsection{The \texttt{"full"} strategy}

The recommended strategy for entity-disjoint settings with
entity-correlated missingness is \texttt{"full"}. It combines a neutral
similarity score (missing $\to 0.5$) with explicit binary indicators:

\begin{example}
res <- mv_compare_fields(xa, xb, method = "full")
res
#> $jw
#> [1] 0.9666667 0.8162963 0.5000000 0.5000000
#>
#> $lv
#> [1] 0.8333333 0.4444444 0.5000000 0.5000000
#>
#> $any_miss
#> [1] 0 0 1 1
#>
#> $both_miss
#> [1] 0 0 0 1
\end{example}

The classifier receives a neutral similarity of~0.5 for all pairs
involving a missing field (neither a positive nor a negative signal for
the match decision based on similarity alone), and the indicator columns
allow the model to learn a negative weight for the \texttt{both\_miss}
indicator---correcting the bias rather than suppressing information.

\subsection{Building a feature matrix}

\texttt{mv\_pair\_features(d, pairs, text\_cols, method, ...)} applies a
comparison strategy across all pairs in a candidate set to produce the
full $n_{\text{pairs}} \times p$ feature matrix used by classifiers:

\begin{example}
feats_full <- mv_pair_features(records, pairs,
                               text_cols = c("name", "city"),
                               method    = "full")
feats_full$X
#>    name_jw  name_lv name_any_miss name_both_miss  city_jw  city_lv
#> 1: 0.966667 0.833333             0              0 0.930556 0.833333
#> 2: 0.500000 0.500000             1              0 0.500000 0.500000
#> 3: 0.500000 0.500000             1              1 0.500000 0.500000
#> 4: 0.500000 0.500000             1              0 0.500000 0.500000
#>    city_any_miss city_both_miss
#> 1:             0              0
#> 2:             1              0
#> 3:             0              0
#> 4:             0              0

feats_full$y  # match labels
#> [1] 1 0 0 0
\end{example}

The function pre-normalizes all text fields via \texttt{norm\_str()}
before comparison, handles both \texttt{NA} and \texttt{""} as missing,
and supports chunk-based processing for large candidate sets via the
\texttt{chunk\_size} argument.

%%  ── 6. Feature-matrix imputation ──────────────────────────────────────
\section{Feature-matrix imputation}
\label{sec:impute}

Some comparison strategies---notably \texttt{"complete\_cases"} and
\texttt{"indicator"}---can produce \texttt{NA} entries in the feature
matrix that must be filled before passing data to a classifier.
Additionally, users who chain \pkg{mvher} with their own similarity
routines may obtain feature matrices with residual \texttt{NA}s from
other sources.

\pkg{mvher} provides seven imputation strategies for feature matrices,
designed around a two-step fit/apply API that prevents leakage:

\begin{example}
Xtr <- data.frame(jw = c(0.9, NA, 0.3, 0.8),
                  lv = c(NA, 0.7, 0.5, NA))
Xte <- data.frame(jw = c(NA, 0.6),
                  lv = c(0.4, NA))

# Step 1: fit imputer on training data only
imp <- mv_fit_imputer(Xtr, strategy = "mean")
imp
#> <mvher_imputer>
#>   strategy : mean
#>   columns  : jw lv
#>   fill jw  : 0.6666667
#>   fill lv  : 0.6000000

# Step 2: apply to both sets using *training* statistics
Xtr_clean <- mv_apply_imputer(Xtr, imp)
Xte_clean <- mv_apply_imputer(Xte, imp)   # uses training means
Xte_clean
#>      jw   lv
#> 1: 0.667 0.40
#> 2: 0.600 0.60
\end{example}

The convenience wrapper \texttt{mv\_impute()} combines both steps:

\begin{example}
res <- mv_impute(Xtr, Xte, strategy = "median")
res$X_test
#>     jw  lv
#> 1: 0.85 0.40
#> 2: 0.60 0.60
\end{example}

Table~\ref{tab:impute} summarizes the seven strategies and their
suitability for different scenarios.

\begin{table}[ht]
\centering
\caption{Feature-matrix imputation strategies in \pkg{mvher}.}
\label{tab:impute}
\smallskip
\begin{tabular}{lp{4.2cm}lp{2.5cm}}
\toprule
\textbf{Strategy} & \textbf{Fill value} & \textbf{MAR-valid}
  & \textbf{Best for} \\
\midrule
\texttt{"mean"}      & Column mean (training)  & Yes & Low--moderate missingness \\
\texttt{"median"}    & Column median (training) & Yes & Skewed distributions \\
\texttt{"zero"}      & 0.0                     & No  & Conservative baselines \\
\texttt{"half"}      & 0.5                     & No  & Neutral probability baseline \\
\texttt{"hot\_deck"} & Random draw from training & Yes & Preserving marginal dist. \\
\texttt{"knn"}       & $k$-NN Euclidean average & Yes & Entity-correlated missingness \\
\texttt{"none"}      & No fill (NA retained)   & ---  & Tree-based classifiers \\
\bottomrule
\end{tabular}
\end{table}

The \texttt{"knn"} strategy is recommended when missingness exceeds
$\approx 30\%$ and is entity-correlated (ICC $> 0.3$). For each row
with missing values in the test set, the imputer locates the $k = 5$
(default) nearest complete rows in the \emph{training} set (by
Euclidean distance on observed columns) and fills missing values with
their average.  Because the reference pool is always the training set,
this is immune to leakage even when applied iteratively.

\subsection{Skipping imputation for native-MV classifiers}
\label{sec:impute:native}

Imputation is not always necessary.  Several widely used tree-based
classifiers handle \texttt{NA}s directly inside the learning algorithm,
without requiring a complete feature matrix.
\CRANpkg{xgboost} \citep{ChenGuestrin2016} and \CRANpkg{lightgbm}
\citep{Ke2017} learn, at each tree node, the branch direction that
minimizes the loss function for missing-valued observations; the
direction is then reused at prediction time.
\CRANpkg{ranger} \citep{WrightZiegler2017} supports a similar
learned-direction mechanism via \texttt{na.action = "na.learn"}.
Classical random forests \citep{Breiman1984} fall back on
surrogate splits---secondary variables that mimic the primary split
when the primary value is missing.

For these classifiers, passing \texttt{strategy = "none"} to
\texttt{mv\_fit\_imputer()} (or \texttt{mv\_impute()}) leaves residual
\texttt{NA}s in place, and the classifier handles them natively.  This
avoids the risk of imputation introducing bias when missingness is
entity-correlated.

Critically, \emph{field-comparison strategy still matters even for
native-MV classifiers}.  If \texttt{"current"} is used at the
comparison stage, both-missing pairs enter the feature matrix with
sim~$= 1$---a numerically complete, but semantically wrong, value that
no native-NA mechanism can correct.  We therefore recommend pairing
native-MV classifiers with \texttt{"indicator"} or \texttt{"full"}:
the \texttt{both\_miss} indicator column gives the tree an explicit
signal to learn a negative weight for both-missing pairs, while
residual \texttt{NA}s from the continuous similarity columns are
handled natively.

\texttt{mv\_recommend()} supports this workflow via the
\texttt{classifier = "tree\_native"} argument:

\begin{example}
rec <- mv_recommend(field_miss_rate = 0.25, entity_corr = 0.6,
                    split_type = "entity",
                    classifier = "tree_native")
print(rec)
#> mvher recommendation
#>   field_method    : full
#>   imputer_strategy: none
#>   rationale:
#>     Native-MV classifier requested (e.g. xgboost, lightgbm, ranger):
#>     imputation skipped (strategy = 'none'). Residual NAs in the
#>     feature matrix are left for the classifier's native
#>     missing-value mechanism. Field-comparison method 'full' still
#>     applied to avoid the both-null=1 artifact.
\end{example}

%%  ── 7. Strategy recommendation ────────────────────────────────────────
\section{Strategy recommendation}
\label{sec:recommend}

\pkg{mvher} provides a strategy table and an automated recommendation
engine to help practitioners select the appropriate combination of
field-comparison method and imputation strategy.

\texttt{mv\_strategy\_table()} returns a 15-row data frame covering all
eight field-comparison methods and seven imputation strategies, with
columns indicating whether each strategy is valid under MCAR, MAR, or
MNAR; whether it adds features to the model; whether it requires
entity labels; and its computational cost.

\texttt{mv\_recommend()} encodes the following decision rules into an
S3 object of class \texttt{mvher\_recommendation}:

\begin{enumerate}
  \item If the field missingness rate is $< 2\%$: use \texttt{"current"}
    and \texttt{"mean"} (near-zero missingness; backward compatibility
    with existing pipelines).
  \item If the split is entity-disjoint and the entity ICC $> 0.3$:
    use \texttt{"full"} (neutral sim~$+ $ indicators); use
    \texttt{"knn"} imputation if missingness $\geq 30\%$, else
    \texttt{"mean"}.
  \item If the field missingness rate is $\geq 15\%$: use
    \texttt{"neutral"} and \texttt{"mean"} (or \texttt{"knn"} if $\geq
    30\%$).
  \item Otherwise: use \texttt{"neutral"} and \texttt{"mean"}.
\end{enumerate}

\begin{example}
# High entity-level correlation under entity-disjoint split
rec <- mv_recommend(field_miss_rate = 0.25, entity_corr = 0.6,
                    split_type = "entity")
print(rec)
#> <mvher_recommendation>
#>   field_method     : full
#>   imputer_strategy : mean
#>   rationale        : Entity-disjoint split with high entity-level
#>                      missingness correlation (ICC=0.60). Use 'full'
#>                      to suppress the both-null artifact and expose
#>                      missingness patterns as explicit features.
\end{example}

%%  ── 8. Case study ─────────────────────────────────────────────────────
\section{End-to-end case study}
\label{sec:casestudy}

We demonstrate the full \pkg{mvher} workflow on a synthetic dataset of
24 records (8 entities $\times$ 3 records each) with 20\% random
name-field missingness and 15\% random city-field missingness.

\begin{example}
set.seed(42)
n <- 24
d <- data.frame(
  id        = seq_len(n),
  entity_id = rep(1:8, each = 3L),
  name      = ifelse(runif(n) < 0.2, NA_character_,
               paste0(sample(c("Smith","Jones","Brown","White"),
                             n, TRUE), " ",
                      sample(c("Alice","Bob","Carol","Dan"),
                             n, TRUE))),
  city      = ifelse(runif(n) < 0.15, "",
               sample(c("London","Paris","Rome"), n, TRUE)),
  stringsAsFactors = FALSE
)
\end{example}

\textbf{Step 1: Diagnose.}

\begin{example}
mv_field_summary(d, c("name", "city"))
#>    column n_total n_empty  pct_miss
#> 1:   name      24       5 0.2083333
#> 2:   city      24       4 0.1666667

mv_entity_correlation(d, c("name", "city"))
#>    column  icc interpretation
#> 1:   name 0.19 low (MCAR-like)
#> 2:   city 0.17 low (MCAR-like)
\end{example}

\textbf{Step 2: Recommend.}

\begin{example}
rec <- mv_recommend(field_miss_rate = c(0.21, 0.17),
                    entity_corr = NULL,
                    split_type  = "entity")
cat(rec$field_method, "/", rec$imputer_strategy, "\n")
#> neutral / mean
\end{example}

\textbf{Step 3: Entity-disjoint 70/30 split.}

\begin{example}
set.seed(1)
ents     <- unique(d$entity_id)
test_ent <- sample(ents, floor(0.3 * length(ents)))
d_tr     <- d[!d$entity_id %in% test_ent, ]
d_te     <- d[ d$entity_id %in% test_ent, ]
\end{example}

\textbf{Step 4: Enumerate all candidate pairs within each split and
extract features.}

\begin{example}
make_all_pairs <- function(ids) {
  ids <- sort(unique(ids)); n <- length(ids)
  if (n < 2) return(data.frame(id1 = integer(), id2 = integer()))
  data.frame(
    id1 = rep(ids[-n], times = rev(seq_len(n - 1))),
    id2 = unlist(lapply(2:n, function(i) ids[i:n]))
  )
}
pairs_tr <- make_all_pairs(d_tr$id)
pairs_te <- make_all_pairs(d_te$id)

ftr <- mv_pair_features(d_tr, pairs_tr, c("name", "city"),
                        method = rec$field_method)
fte <- mv_pair_features(d_te, pairs_te, c("name", "city"),
                        method = rec$field_method)
\end{example}

\textbf{Step 5: Impute feature matrices.}

\begin{example}
res <- mv_impute(ftr$X, fte$X, strategy = rec$imputer_strategy)

# Training set: no NAs remain
anyNA(res$X_train)
#> [1] FALSE

# Test set: filled using training-set means
anyNA(res$X_test)
#> [1] FALSE

head(res$X_train, 3)
#>    name_jw  name_lv  city_jw  city_lv
#> 1: 0.500000 0.500000 0.944444 0.857143
#> 2: 0.500000 0.500000 0.944444 0.857143
#> 3: 0.888889 0.666667 0.944444 0.857143
\end{example}

\textbf{Step 6: Train a classifier and evaluate.}

The resulting matrices \texttt{res\$X\_train} and
\texttt{res\$X\_test}, together with the match labels \texttt{ftr\$y}
and \texttt{fte\$y}, can be passed directly to any binary classifier in
R.  For example, using logistic regression:

\begin{example}
train_df   <- cbind(res$X_train, match = ftr$y)
test_df    <- cbind(res$X_test,  match = fte$y)

fit <- glm(match ~ ., data = train_df, family = binomial)
pred <- predict(fit, newdata = test_df, type = "response")

# AUC (requires pROC)
# pROC::auc(fte$y, pred)
\end{example}

This completes the end-to-end pipeline from raw records with missing
values to a calibrated classifier output.

%%  ── 9. Comparison with related packages ───────────────────────────────
\section{Comparison with related packages}
\label{sec:compare}

Several R packages support entity resolution or missing-value handling,
but none addresses the intersection targeted by \pkg{mvher}.

\CRANpkg{RecordLinkage} \citep{Sariyar2010} provides blocking,
comparison, and classification tools for ER, including the
\texttt{jarowinkler} and \texttt{levenshtein} comparison functions.
However, missing values are handled by \texttt{NA} propagation: any
pair involving a missing field receives \texttt{NA} for that field's
similarity, which is then filled by the user (typically with zero or
dropped entirely). No ER-specific strategy is offered.

\CRANpkg{fastLink} \citep{Enamorado2019} implements the Fellegi-Sunter
model with EM estimation of match and non-match parameters. It encodes a
three-level comparison that partially accommodates missingness, but the
gamma encoding is directly tied to the Fellegi-Sunter likelihood and
is not easily adapted to machine-learning classifiers.

\CRANpkg{mice} \citep{vanBuuren2018} and \CRANpkg{missForest}
\citep{Stekhoven2012} provide state-of-the-art multiple imputation for
general tabular data. They operate on wide record tables, not on
pairwise feature matrices, and neither is aware of the ER-specific
both-null artifact or the leakage risk under entity-disjoint splits.

\pkg{mvher} complements these packages by operating specifically at the
field-comparison and feature-extraction stages of the ER pipeline.  It
can be combined with blocking from \CRANpkg{RecordLinkage} or
\pkg{GCMER} \citep{Degras2025} and classification from any standard
R package (e.g., \CRANpkg{glmnet}, \CRANpkg{ranger}).

Table~\ref{tab:compare} summarizes the key distinctions.

\begin{table}[ht]
\centering
\caption{Comparison of R packages for entity resolution and
  missing-value handling.}
\label{tab:compare}
\smallskip
\begin{tabular}{lcccc}
\toprule
\textbf{Package} & \textbf{ER pipeline} & \textbf{MV at comparison}
  & \textbf{Leakage-safe API} & \textbf{Both-null fix} \\
\midrule
\pkg{RecordLinkage} & Yes & Partial (NA propagation) & No  & No  \\
\pkg{fastLink}      & Yes & Partial (FS only)         & No  & Partial \\
\pkg{mice}          & No  & No (record-level)        & No  & No  \\
\pkg{mvher}         & No  & Yes (8 strategies)       & Yes & Yes \\
\bottomrule
\end{tabular}
\end{table}

%%  ── 10. Summary ───────────────────────────────────────────────────────
\section{Summary}
\label{sec:summary}

We have presented \pkg{mvher}, an R package for principled missing-value
handling in entity resolution pipelines. The package addresses a
concrete, often overlooked failure mode: the spurious both-null-$=1$
similarity score produced by standard string metrics, which inflates
estimated match probabilities and creates covariate shift under
entity-disjoint evaluation. The \texttt{"full"} comparison strategy
(neutral similarity plus binary missingness indicators), combined with
the leakage-safe fit/apply imputation API, provides a practical
remedy that integrates cleanly with any downstream classifier.

The key contributions of \pkg{mvher} are:

\begin{enumerate}
  \item A unified, ER-aware framework for diagnosing, quantifying, and
    handling missing values at the field-comparison stage.
  \item An eight-strategy field-comparison API that makes the treatment
    of missing values an explicit, documented modeling choice rather
    than an implicit default.
  \item A seven-strategy feature-matrix imputation API with a
    fit-on-training/apply-to-test design that prevents data leakage
    under entity-disjoint splits.
  \item An automated recommendation engine that maps data
    characteristics to strategy combinations, lowering the
    barrier for practitioners unfamiliar with the theoretical
    foundations.
\end{enumerate}

Future work will extend the package in several directions.  First, the
current diagnosis functions are heuristic; a more principled test for
entity-correlated missingness under the ER data-generating model would
improve the accuracy of the recommendation engine.  Second, the
imputation module currently handles feature-matrix NAs only; incorporating
record-level imputation of raw fields before comparison would allow a
more comprehensive two-stage approach.  Third, benchmarks on real-world
ER datasets with controlled missingness rates would quantify the
practical benefit of the \texttt{"full"} strategy relative to the
\texttt{"current"} default across a range of classifiers and missingness
levels.

%%  ── Acknowledgments ───────────────────────────────────────────────────
%\section*{Acknowledgments}



%%  ── Bibliography ──────────────────────────────────────────────────────
\bibliographystyle{abbrvnat}
\bibliography{mvher}

\end{article}
\end{document}
